\documentclass[a4paper,12pt]{article}

\usepackage[T2A]{fontenc}
\usepackage[utf8]{inputenc}
\usepackage[english,russian]{babel}

\usepackage{amsmath}
\usepackage{amssymb}
\usepackage[margin=2cm]{geometry}
\usepackage{parskip}
\usepackage{enumitem}
\usepackage{booktabs}
\usepackage{float}
\usepackage{tikz}
\usetikzlibrary{shapes.multipart,positioning,arrows.meta,calc}
\usepackage[hidelinks]{hyperref}

\setlist{itemsep=0.2em,topsep=0.35em}
\sloppy

\tikzset{
  cls/.style={
    rectangle split,
    rectangle split parts=3,
    rectangle split part align={center,left,left},
    draw,
    font=\scriptsize,
    align=left,
    inner xsep=4pt,
    inner ysep=2.5pt,
    text width=4.15cm
  },
  smallcls/.style={
    rectangle split,
    rectangle split parts=2,
    rectangle split part align={center,left},
    draw,
    font=\scriptsize,
    align=left,
    inner xsep=4pt,
    inner ysep=2.5pt,
    text width=3.5cm
  },
  obj/.style={
    draw,
    font=\scriptsize,
    align=left,
    inner sep=3.5pt,
    text width=3.35cm
  },
  inherit/.style={{Triangle[open,length=7pt,width=7pt]}-,semithick},
  comp/.style={{Diamond[fill,length=5.5pt,width=5.5pt]}-,semithick},
  assoc/.style={-{Stealth[length=6pt]},semithick},
  depend/.style={-{Stealth[length=6pt]},dashed,semithick},
  rel/.style={font=\scriptsize,fill=white,inner sep=1.5pt}
}

\title{Моделирование обслуживания в филиале банка}
\author{Кирилл Аверченко, Анжелика Гаврилова}
\date{October 2026}

\begin{document}

\begin{center}
{\Large\bfseries Моделирование обслуживания в филиале банка}\\[0.55em]
Отчёт по заданию\\[0.25em]
\end{center}

\section{Уточнённая постановка задачи}

Нужно промоделировать обслуживание клиентов в одном филиале банка и по результатам месяца понять, как прибыль и потери зависят от числа клерков.

В филиале \(N\) клерков, \(2 \le N \le 7\), и одна общая очередь длины не больше \(K\), \(10 \le K \le 25\). Заявки в течение дня нумеруются с~1. На табло показывается номер клиента, только что взятого на обслуживание, и номер окна.

Филиал работает четыре недели. Пять дней открыты с 9:00 до 17:00, суббота~--- с 10:00 до 16:00, воскресенье закрыто. Четыре недели взяты вместо календарного месяца, чтобы каждый день недели встретился одинаково часто и сравнение числа клерков не зависело от того, на какой день выпало первое число. Обед длится 30 минут. Окна обеда сдвинуты и не пересекаются: полоса обедов стоит по центру рабочего дня, клерк \(i\) начинает обед на \(i\) интервалов позже первого. Если в номинальное время клерк занят, он сначала дообслуживает клиента и только потом уходит. На обеде из очереди он никого не берёт.

Интервал между клиентами и длительность обслуживания задаются отдельно. Каждая величина либо равномерна на отрезке \([\min,\max]\), либо нормальна с параметрами \(\mu\) и \(\sigma\) и отсекается в этот же отрезок. Нулевой интервал между приходами заменяется одной миллисекундой, иначе при нулевом диапазоне модель не продвигается по времени. Прибыль с обслуженного клиента равномерна на целых рублях от 50 до 3000. Дневная зарплата клерка~--- 2000 рублей, она платится за каждый из 24 рабочих дней, включая субботу.

Плотность потока не постоянна. Базовый интервал \(\Delta\)~--- это интенсивность в среднем по открытому времени при очереди не длиннее 7 человек. Фактический интервал равен \(\Delta / m\), где
\[
m = m_{\text{нед}}(d)\, m_{\text{час}}(h)\, m_{\text{оч}}(L).
\]
День \(d = 0,\ldots,5\) (понедельник\ldots суббота), \(l_d\)~--- длина дня в часах, за неделю 46 рабочих часов. Веса \((d+1)l_d\) в сумме дают 156, поэтому
\[
m_{\text{нед}}(d) = \frac{46}{156}(d+1)
\]
имеет среднее 1 по рабочим часам: понедельник около \(0{,}29\), суббота около \(1{,}77\). Час внутри дня \(h = 0,\ldots,H-1\), \(H = 8\) или \(H = 6\):
\[
m_{\text{час}}(h) = \frac{2(h+1)}{H+1}.
\]
Среднее по часам дня тоже равно 1. Пока в очереди не больше 7 человек, \(m_{\text{оч}} = 1\). Дальше поток редеет, но не обрывается:
\[
m_{\text{оч}}(L) = \frac{K-L+1}{K-6}, \qquad L > 7.
\]
При \(L = K\) множитель положителен, поэтому отказ из-за полной очереди по-прежнему возможен. Такой клиент номер не получает и в среднее ожидание не входит.

\(\Delta\) разыгрывается один раз и хранится как остаток работы. Работа убывает со скоростью \(m\). На границе часа и при изменении длины очереди остаток пересчитывается, а уже назначенное время прихода сдвигается. Заново случайное число не берётся: повторный розыгрыш обрезал бы длинные интервалы.

В момент закрытия ожидающий приход снимается. Все, кто стоит в очереди, уходят и считаются потерянными; в их ожидание пишется время от прихода до закрытия. Клиента у окна клерк дообслуживает и прибыль по нему получает, но следующего из очереди уже не вызывает. Следующий интервал разыгрывается в момент следующего открытия.

Шаг моделирования~--- от 10 до 60 минут. Внутри шага модель событийная: обрабатываются все события не позже конца шага, и только потом интерфейс показывает кадр. Пока филиал закрыт, один шаг переходит к ближайшему событию, обычно к следующему открытию. Ночь и воскресенье в средние по времени не входят.

Считаются число обслуженных и потерянных клиентов, отдельно отказы из-за полной очереди и уход при закрытии, минимальная, средняя и максимальная длина очереди, среднее ожидание, средняя занятость клерков и прибыль. Средняя длина очереди и занятость~--- средние по открытому времени: между событиями состояние постоянно, поэтому к сумме добавляется длина, умноженная на длину промежутка. Занятым считается только окно, которое обслуживает клиента. Прибыль равна сумме прибылей обслуженных клиентов минус зарплата.

Цель прогонов~--- увидеть, как с ростом \(N\) меняется прибыль и где филиал теряет клиентов: на полной очереди, на закрытии или на простое уже нанятых клерков.

\newpage

\section{Диаграмма классов}

На рисунке~\ref{fig:classes} показаны классы, которые образуют систему. Вложенные в \texttt{Bank} типы \texttt{Timed}, \texttt{Waiting}, \texttt{Clerk} и \texttt{Plan} вынесены в спецификацию: на диаграмме они сделали бы схему нечитаемой, а наружу из \texttt{Bank} не видны. \texttt{ClerkView}~--- это фаза клерка внутри снимка, поэтому нарисован рядом с \texttt{BankSnapshot}.

\begin{figure}[H]
\centering
\begin{tikzpicture}
  \node[cls] (mw) at (0,0) {%
    \textbf{MainWindow}
    \nodepart{second}
    bank : Bank\\
    view : BankWidget\\
    log, stats\\
    поля \(N\), \(K\), шага,\\
    законов и прибыли
    \nodepart{third}
    + startOver()\\
    + stepOnce()\\
    + togglePlay()\\
    + runToEnd()\\
    -- readRange(...)
  };
  \node[cls] (bw) at (5.7,0) {%
    \textbf{BankWidget}
    \nodepart{second}
    snap : BankSnapshot\\
    hasSnap : bool
    \nodepart{third}
    + setSnapshot(s)\\
    \# paintEvent()
  };
  \node[cls] (bank) at (0,-6.4) {%
    \textbf{Bank}
    \nodepart{second}
    params : BankParams\\
    gen : Generators\\
    events : set\\
    queue, clerks, plan\\
    суммы статистики
    \nodepart{third}
    + advanceTo(t)\\
    + advanceStep(dt)\\
    + snapshot()\\
    -- speed(), retarget()
  };
  \node[smallcls] (snap) at (5.7,-6.4) {%
    \textbf{BankSnapshot}
    \nodepart{second}
    time, open, done\\
    queue : vector<int>\\
    clerks : vector\\
    served, lost, profit\\
    queueAvg, occupancy
  };
  \node[smallcls] (cv) at (11.1,-6.4) {%
    \textbf{ClerkView}
    \nodepart{second}
    phase : Free | Serving | Lunch\\
    client : int\\
    until : время конца фазы
  };
  \node[smallcls] (gen) at (0,-11.2) {%
    \textbf{Generators}
    \nodepart{second}
    rnd : mt19937\\
    + randomInt(a, b)\\
    + sample(...)
  };
  \node[smallcls] (par) at (5.7,-11.2) {%
    \textbf{BankParams}
    \nodepart{second}
    clerks, limit : int\\
    arrival, service : DistRange\\
    profitMin, profitMax : int\\
    seed : long long
  };
  \node[smallcls] (dist) at (11.1,-11.2) {%
    \textbf{DistRange}
    \nodepart{second}
    normal : bool\\
    minMs, maxMs : int\\
    meanMs, sigmaMs : double
  };

  \draw[assoc] (mw.east) -- node[rel,above]{показывает} (bw.west);
  \draw[assoc] (mw.south) -- node[rel,right]{1} (bank.north);
  \draw[depend] (bw.south) -- (snap.north);
  \draw[comp] (snap.east) -- node[rel,above]{\(N\)} (cv.west);
  \draw[comp] (bank.south) -- node[rel,right]{1} (gen.north);
  \draw[comp] (bank.south east) -- ++(0,-0.45) -| node[rel,above,pos=0.25]{1} (par.north);
  \draw[comp] (par.east) -- node[rel,above]{2} (dist.west);
\end{tikzpicture}
\caption{Диаграмма классов. Сплошная стрелка~--- использование, пунктир~--- зависимость от снимка, закрашенный ромб~--- владение.}
\label{fig:classes}
\end{figure}

\texttt{MainWindow} владеет единственным объектом \texttt{Bank} и создаёт его заново при каждом запуске. \texttt{BankWidget} модель не вызывает: он получает уже собранный \texttt{BankSnapshot}. Генератор и параметры живут внутри модели, интерфейс их не меняет посреди прогона.

\section{Текстовые спецификации основных классов}

\subsection*{MainWindow}
Главное окно. Строит форму параметров, зал и панель статистики, хранит \texttt{unique\_ptr} на текущую модель и таймер автоматического шага.

Операции \texttt{startOver}, \texttt{stepOnce}, \texttt{togglePlay} и \texttt{runToEnd} вызываются кнопками. \texttt{startOver} проверяет диапазоны, берёт новый сид, создаёт \texttt{Bank} и сразу открывает утро понедельника. \texttt{stepOnce} двигает модель на шаг из поля <<Шаг>>. \texttt{runToEnd} останавливает таймер и вызывает \texttt{advanceTo} до конца четырёх недель, без промежуточной перерисовки. \texttt{readRange} переводит минуты из полей в миллисекунды и выбирает равномерный или нормальный закон. \texttt{refresh} забирает снимок и обновляет зал, числа и журнал.

Связи: один \texttt{Bank}, один \texttt{BankWidget}, набор виджетов Qt для ввода.

\subsection*{BankWidget}
Виджет зала. Хранит последний снимок и рисует его в \texttt{paintEvent}: часы и признак <<открыто>> или <<закрыто>>, табло, окна с фазой и номером клиента, очередь кружками с номерами.

Операция \texttt{setSnapshot} заменяет снимок и заказывает перерисовку. К модели виджет обращений не имеет.

\subsection*{Bank}
Имитационная модель. Время~--- целые миллисекунды от полуночи первого понедельника. Календарь событий~--- \texttt{std::set} структур \texttt{Timed} с ключом (время, приоритет, порядковый номер). Приоритет задан порядком перечисления: закрытие, открытие, конец обслуживания, конец обеда, начало обеда, приход, граница часа. Порядковый номер делает одинаковые времена однозначными, поэтому один и тот же сид повторяет прогон.

Состояние: параметры, генератор, очередь \texttt{Waiting}, вектор \texttt{Clerk} длины \(N\), план одного будущего прихода, номерки текущего дня, табло, накопители статистики и журнал. Длина очереди \(L\) хранится отдельно, чтобы не пересчитывать её на каждом событии.

Открытые операции:
\begin{itemize}
  \item \texttt{advanceTo(t)} обрабатывает все события не позже \(t\), но не дальше конца месяца. Если событий больше нет, часы всё равно доводятся до \(t\), и за открытый кусок дописывается статистика.
  \item \texttt{advanceStep(dt)} в часы работы равносилен шагу длины \(dt\). Когда филиал закрыт, переходит к ближайшему событию, а если событий не осталось~--- к концу месяца.
  \item \texttt{snapshot} копирует очередь, окна и статистику в структуру без ссылок на внутренние контейнеры.
  \item \texttt{now}, \texttt{endTime}, \texttt{finished} нужны окну, чтобы понять, есть ли ещё что считать.
\end{itemize}

Закрытые операции разбирают события. \texttt{speed} считает \(m\) по дню, часу и \(L\). \texttt{armNewInterval} разыгрывает новый \(\Delta\) и ставит приход. \texttt{retarget} списывает уже потраченную работу и переставляет приход при новой скорости. \texttt{tryAssign} после освобождения окна либо отправляет клерка на обед, либо берёт первого из очереди, либо оставляет окно свободным. На закрытии \texttt{cancelArrival} снимает приход, очередь списывается в потери, зарплата прибавляется один раз за день.

Инвариант открытого филиала: нет одновременно свободного окна и непустой очереди. Свободное окно сразу забирает первого ожидающего, если обед этому клерку ещё не начался.

\subsection*{Generators}
Источник случайных чисел, один на модель. Внутри \texttt{std::mt19937}, инициализируется сидом прогона.

\texttt{randomInt(a, b)} даёт целое на отрезке с включёнными концами. \texttt{randomDouble(a, b, k)} даёт число с \(k\) знаками после запятой; модель им не пользуется, метод оставлен как общий датчик. \texttt{sample} при равномерном законе сводится к \texttt{randomInt}. При нормальном законе значение из \(\mathcal N(\mu,\sigma)\) отбрасывается, пока не попадёт в \([\min,\max]\), но не дольше 10\,000 попыток; иначе берётся \(\mu\), обрезанное границами.

\subsection*{BankParams и DistRange}
\texttt{BankParams}~--- значение, с которым создаётся модель: \(N\), \(K\), два закона, границы прибыли, сид. После создания модель эти поля не меняет. Шаг моделирования в параметры не входит: шаг принадлежит окну и может быть другим на следующем нажатии.

\texttt{DistRange} описывает один закон. Флаг \texttt{normal} выбирает распределение. Границы хранятся в миллисекундах, \(\mu\) и \(\sigma\) тоже переведены в миллисекунды ещё в окне.

\subsection*{BankSnapshot и ClerkView}
\texttt{BankSnapshot}~--- кадр для интерфейса и для отчёта по эксперименту. В нём часы, открыт ли филиал, закончен ли месяц, выданные сегодня номера, табло, номера очереди, вид каждого клерка, счётчики потерь по двум причинам, минимумы и средние, выручка, зарплата и прибыль, а также хвост журнала.

\texttt{ClerkView} содержит только то, что нужно нарисовать: фазу, номер клиента и абсолютное время конца обслуживания или обеда. Деньги и момент прихода клиента в снимок клерка не входят.

\subsection*{Вложенные типы Bank}
\texttt{Timed}~--- одно будущее событие. \texttt{Waiting}~--- человек в очереди: дневной номер и время прихода. \texttt{Clerk} расширяет фазу из \texttt{ClerkView} данными текущего обслуживания: когда клиент пришёл, когда его взяли, длительность, прибыль, номинальное начало обеда и признак, что обед уже начат. \texttt{Plan} хранит единственный будущий приход: остаток работы, скорость, с которой этот остаток списывался, и ключ события в календаре, чтобы его можно было удалить и вставить снова.

\section{Диаграмма объектов}

На рисунке~\ref{fig:objects} разобран один момент работы: понедельник первой недели, 12:30, \(N = 3\), \(K = 15\). При таких \(N\) обед первого окна длится с 12:15 до 12:45, второго~--- с 12:45 до 13:15, третьего~--- с 13:15 до 13:45. Значения атрибутов подобраны так, чтобы выполнялся инвариант модели; это схема связей, а не протокол конкретного сида.

\begin{figure}[H]
\centering
\begin{tikzpicture}
  \node[obj] (bank) at (0,0) {%
    \underline{branch : Bank}\\
    cur = 12:30\\
    open = true\\
    \(L = 2\)\\
    выдано номеров = 7\\
    board = клиент 5, окно 3
  };
  \node[obj] (gen) at (4.8,0) {%
    \underline{rnd : Generators}\\
    сид прогона\\
    уже использовано\\
    на номера 1--7
  };
  \node[obj] (plan) at (9.6,0) {%
    \underline{next : Plan}\\
    active = true\\
    at = 12:36\\
    work~--- остаток \(\Delta\)\\
    \(m\)~--- скорость этого часа
  };
  \node[obj] (c0) at (0,-3.7) {%
    \underline{c0 : Clerk}\\
    phase = Lunch\\
    client = 0\\
    until = 12:45\\
    lunchDone = true
  };
  \node[obj] (c1) at (4.8,-3.7) {%
    \underline{c1 : Clerk}\\
    phase = Serving\\
    client = 4\\
    arrived = 12:02\\
    started = 12:14\\
    until = 12:38
  };
  \node[obj] (c2) at (9.6,-3.7) {%
    \underline{c2 : Clerk}\\
    phase = Serving\\
    client = 5\\
    arrived = 12:18\\
    started = 12:25\\
    until = 12:51
  };
  \node[obj] (w6) at (4.8,-7.15) {%
    \underline{w6 : Waiting}\\
    number = 6\\
    arrival = 12:22
  };
  \node[obj] (w7) at (9.6,-7.15) {%
    \underline{w7 : Waiting}\\
    number = 7\\
    arrival = 12:29
  };

  \draw[comp] (bank.east) -- node[rel,above]{1} (gen.west);
  \draw[comp] (bank.north) -- ++(0,0.55) -| node[rel,above,pos=0.25]{0..1} (plan.north);
  \draw[comp] (bank.south) -- ++(0,-0.4) coordinate (hub);
  \draw (hub) -- (c0.north);
  \draw (hub) -| (c1.north);
  \draw (hub) -| (c2.north);
  \draw[comp] (bank.west) -- ++(-0.55,0) |- node[rel,pos=0.25,left]{0..\(K\)} (w6.west);
  \draw (w6.east) -- (w7.west);
\end{tikzpicture}
\caption{Диаграмма объектов в понедельник в 12:30. Окно~1 на обеде, окна~2 и~3 заняты, поэтому клиенты 6 и~7 стоят в очереди. Ближайший приход назначен на 12:36.}
\label{fig:objects}
\end{figure}

В календаре \texttt{branch} в этот момент лежат, среди прочих, конец обслуживания клиента~4 в 12:38, конец обеда окна~1 в 12:45, приход в 12:36 и граница часа в 13:00. Приход стоит в календаре позже конца обслуживания, хотя на часах он раньше: порядок в наборе задаётся временем, а не порядком вставки. Когда в 12:36 придёт следующий клиент, оба рабочих окна ещё заняты, и он встанет в очередь третьим, если к этому времени никого не вызовут. В 12:38 окно~2 освободится и заберёт клиента~6; план прихода при этом будет пересчитан, потому что длина очереди уменьшится.

\section{Инструментальные средства}

Программа написана на C++17. Сборка описана в \texttt{qmake}-файле \texttt{Model.pro} и выполняется комплектом Desktop Qt 5.12.12 MinGW 64-bit в Qt Creator. Компилятор комплекта~--- MinGW GCC 7.3. Исходники читаются как UTF-8: иначе русские строки в интерфейсе на этой связке портятся.

Из стандартной библиотеки используются \texttt{std::mt19937} и распределения из \texttt{<random>}, \texttt{std::set} как календарь событий, \texttt{std::deque} как очередь и журнал, \texttt{std::vector} в снимке. Интерфейс сделан на Qt Widgets: главное окно, формы, счётчики, таймер шага и собственный виджет с переопределённым \texttt{paintEvent}. Отдельной графической библиотеки нет. Модель под Qt не подключается и собирается без него: окно только создаёт \texttt{Bank} и читает снимок.

\newpage

\section{Файловая структура}

Каталог проекта~--- \texttt{Documents/Model}. Рядом Qt Creator кладёт теневую сборку \texttt{build-Model-Desktop\_Qt\_5\_12\_12\_MinGW\_64\_bit-Debug}; в отчёт она не входит.

\begin{figure}[H]
\centering
\begin{tikzpicture}[
  every node/.style={font=\ttfamily\small,inner sep=2pt},
  dir/.style={font=\ttfamily\bfseries\small}
]
  \node[dir] (root) {Model};
  \node[dir, below left=4mm and 28mm of root] (h) {заголовки};
  \node[dir, below right=4mm and 28mm of root] (s) {реализация};
  \node[below=2mm of h, anchor=north] {%
    \begin{tabular}{@{}l@{}}
      bank.h\\
      bankwidget.h\\
      generators.h\\
      mainwindow.h
    \end{tabular}};
  \node[below=2mm of s, anchor=north] {%
    \begin{tabular}{@{}l@{}}
      bank.cpp\\
      bankwidget.cpp\\
      generators.cpp\\
      main.cpp\\
      mainwindow.cpp\\
      Model.pro
    \end{tabular}};
  \draw (root.south) -- ++(0,-2mm) -| (h.north);
  \draw (root.south) -- ++(0,-2mm) -| (s.north);
\end{tikzpicture}
\caption{Файлы программной системы.}
\label{fig:files}
\end{figure}

\texttt{bank.h} и \texttt{bank.cpp} содержат модель, снимок и параметры. \texttt{generators.h} и \texttt{generators.cpp}~--- датчик. \texttt{bankwidget.h} и \texttt{bankwidget.cpp}~--- рисование зала. \texttt{mainwindow.h} и \texttt{mainwindow.cpp}~--- форма, кнопки и вывод статистики. \texttt{main.cpp} создаёт приложение Qt и окно. \texttt{Model.pro} перечисляет эти файлы, включает модуль Widgets, стандарт C++17 и чтение исходников в UTF-8.

\section{Пользовательский интерфейс}

Окно делится на три части. Слева параметры: число клерков от 2 до 7, предел очереди от 10 до 25, шаг от 10 до 60 минут, закон и границы интервала между клиентами, закон и границы обслуживания, прибыль от и до. Для нормального закона дополнительно задаются \(\mu\) и \(\sigma\) в минутах; при равномерном законе эти поля недоступны.

Кнопка <<Запустить сначала>> проверяет, что минимум не больше максимума и что верх интервала больше нуля, затем берёт новый сид и открывает филиал в понедельник в 9:00. <<Шаг>> продвигает модель на выбранные минуты, а в закрытые часы~--- до ближайшего события. <<Пуск>> повторяет шаг по таймеру, повторное нажатие ставит паузу. <<До конца>> доигрывает месяц без промежуточных кадров: зал и статистика обновляются один раз, когда курсор ожидания снимается.

В центре~--- зал. Сверху часы, день недели, номер недели и подпись, открыт филиал или нет. Под ними табло. Ниже ряд окон: зелёное свободно, синее обслуживает и показывает номер клиента, оранжевое на обеде. Под окнами нарисована очередь, каждый человек~--- кружок с дневным номером, и подпись, сколько номеров уже выдано сегодня.

Справа статистика за уже прожитое открытое время и журнал. В статистике обслуженные, потерянные с разбивкой на отказ и закрытие, минимум, среднее и максимум очереди, среднее ожидание в минутах, занятость в процентах, выручка, зарплата и прибыль. Журнал хранит последние 400 строк: открытие, обед, приход, отказ, вызов из очереди, конец обслуживания, закрытие. Статистика при этом копится за весь месяц, а не за эти 400 строк.

\section{Проведённые эксперименты}

Сравнивались пять значений \(N\) при остальных параметрах, совпадающих со значениями окна по умолчанию. Интервал между клиентами равномерный от 0 до 10 минут, обслуживание равномерное от 2 до 30 минут, прибыль от 50 до 3000 рублей, \(K = 15\), сид 7. Каждый прогон~--- четыре недели, то есть 24 рабочих дня. Зарплата поэтому равна \(N \cdot 2000 \cdot 24\).

\begin{table}[H]
\centering
\small
\begin{tabular}{@{}rrrrrrrr@{}}
\toprule
\(N\) & обслужено & отказ & закрытие & очередь, средн. & ожидание, мин & занятость & прибыль, руб. \\
\midrule
2 & 1024 & 11  & 265 & 5,16 & 44,2 & 71,0\% & 1\,494\,465 \\
3 & 1370 & 0   & 214 & 4,01 & 28,0 & 64,7\% & 1\,934\,830 \\
4 & 1597 & 1   & 176 & 2,84 & 17,7 & 57,0\% & 2\,212\,704 \\
5 & 1813 & 0   & 134 & 1,92 & 10,9 & 51,5\% & 2\,509\,451 \\
7 & 2078 & 0   & 82  & 0,67 & 3,4  & 42,0\% & 2\,811\,615 \\
\bottomrule
\end{tabular}
\caption{Зависимость итогов месяца от числа клерков. Сид 7, \(K = 15\), базовые диапазоны.}
\label{tab:clerks}
\end{table}

Минимум очереди во всех пяти прогонах равен 0: филиал каждый день открывается пустым, и этот ноль попадает в минимум. Максимум при \(N\) от 2 до 5 равен 15, то есть очередь хотя бы раз упирается в \(K\). При семи клерках максимум равен 12, предел за месяц не достигается.

Отказ из-за полной очереди почти не влияет на итог: за месяц таких клиентов 11, 0, 1, 0 и 0. Основные потери~--- роспуск очереди в конце дня. С двумя клерками так уходит 265 человек, с семью~--- 82. Среднее ожидание падает с 44 минут до 3,4, средняя длина очереди~--- с 5,2 до 0,7. Занятость при этом снижается с 71\% до 42\%: лишние окна часть дня пустуют, и зарплата растёт на 48\,000 рублей с каждым новым клерком.

Прибыль всё равно растёт вплоть до \(N = 7\), от 1,49 до 2,81 миллиона рублей. Дополнительная выручка от тех, кого успевают обслужить до закрытия, больше дополнительной зарплаты. На этих диапазонах узкое место~--- не потолок \(K\) и не простой, а хвост очереди к закрытию. Простой становится заметен по занятости, но в деньгах ещё не перекрывает потери вечера.

Отдельно проверены два края, чтобы отличить поведение модели от свойств базового потока.

Редкий поток: интервал 30--40 минут, обслуживание 2--5 минут, \(N = 4\), \(K = 12\), сид 1. Обслужено 310 человек, потерь нет, очередь за месяц так и не появилась, ожидание нулевое, занятость 2,4\%. Прибыль 275\,655 рублей при зарплате 192\,000: выручка 467\,655 рублей едва перекрывает оклад. Здесь узкое место уже противоположное~--- простой.

Плотный поток: интервал 5--15 секунд, обслуживание 10--20 минут, \(N = 2\), \(K = 10\), сид 2. Очередь держится на пределе: максимум равен 10, отказов 15\,610. На каждом из 24 закрытий распускаются все 10 ожидающих, это 240 человек. Занятость 93,1\%. Оставшиеся до ста процентов почти целиком приходятся на обязательный получасовой обед каждого клерка, а не на простой. Обслужено только 1400 человек: два окна физически не успевают за таким потоком, и модель в этом месте упирается в \(K\), как и требует постановка.

Третья проверка касается обеда, а не статистики месяца. При \(N = 3\) и интервале в три часа, чтобы утро прошло без клиентов, в 12:15 на обеде ровно первое окно, два других свободны. В 12:45 первое окно уже свободно, на обед уходит второе. Номинальные получасовые окна не накладываются.

По повторным прогонам с одним сидом обслуженные, потери и прибыль совпадают. Это ожидаемо: календарь упорядочен полностью, а генератор инициализируется тем же сидом. В окне каждое <<Запустить сначала>> берёт новый сид, поэтому два нажатия <<До конца>> подряд без перезапуска дают одни и те же числа, а два перезапуска~--- разные.

\end{document}
